\uuid{cwGt}
\exo7id{2725}
\titre{exo7 2725}
\auteur{matexo1}
\organisation{exo7}
\datecreate{2002-02-01}
\isIndication{false}
\isCorrection{false}
\chapitre{Série numérique}
\sousChapitre{Série à  termes positifs}
\cpge{mpsi-mp2i.procedes-sommatoires-discrets.a, pcsi.series-numeriques.a, ptsi.series-numeriques.a, mpsi-mp2i.procedes-sommatoires-discrets.b, pcsi.series-numeriques.b, ptsi.series-numeriques.b, mpsi-mp2i.procedes-sommatoires-discrets.c, pcsi.series-numeriques.c, ptsi.series-numeriques.c, bcpst1.methodes-calcul, mp-mpi.series-numeriques-vectorielles.a, pt.series-numeriques.b, mp-mpi.series-numeriques-vectorielles.b, pc.series-numeriques, psi.series-numeriques.a, pt.series-numeriques.a, bcpst2.series-reelles}
\module{Analyse}
\niveau{L1}
\difficulte{}

\contenu{
\question{
Justifier la convergence et calculer les sommes des s\'eries suivantes 
$$\begin{array}{ll}
\displaystyle \sum_{n\geq 1} \frac 1{n(n+k)} \quad(k\in\N^*), &
\displaystyle\sum \frac 9{(3n+1)(3n+4)}, \\
\displaystyle\sum \frac{n^2+n-3}{n!}, &
\displaystyle\sum_{n\geq 2} \ln \frac{n^2}{n^2-1}.
\end{array}$$
}
}
